% ---- Modele TP PC V2 -----------
\documentclass[a4paper,10pt,french]{scrartcl}
\usepackage[T1]{fontenc}%---pour l'encodage d'entr\'ee
\usepackage{array}%---pour les tableaux
\usepackage{titlesec}%---modification des titres
\usepackage{ulem}%---pour le soulignage
\usepackage[dvipsnames]{xcolor}%---pour avoir d'autres couleurs
\usepackage{color}%---pour avoir les couleurs de base
\usepackage{babel}%---pour les trucs de langue et la francisation

\usepackage{graphicx}%----pour mettre des images
\usepackage[utf8]{inputenc}%---encodage
\usepackage{geometry}%---pour modifier les tailles et mettre a4paper
%\usepackage{awesomebox}%---pour les boites d'exercices, de pbq et de croquis ---d\'esactiv\'e pour les TP de PC
\usepackage{tikz}%---pour dessiner + d\'ependance de chemfig
%\usepackage{tabularx}%---pour dimensionner automatiquement les tableaux avec variable X
\usepackage{fancyhdr}%---pour les en-t\^ete personnalis\'ees
\usepackage{blindtext}%---pour les liens
\usepackage{hyperref}%---pour les liens (\`a mettre en dernier)
\usepackage{caption}%---pour la francisation de la l\'egende table vers Tableau
 \usepackage{float}% --- pour placer les figures et tableaux où l'on souhaite avec argument H
 \usepackage{multicol}
 \usepackage{amsmath}
 %\usepackage{geometry}
 %\geometry{vmargin=2.5cm}

\captionsetup{labelfont=sc}%---pour la francisation de la l\'egende table vers Tableau
\def\frenchtablename{Tableau}%---pour la francisation de la l\'egende table vers Tableau
\pagestyle{fancy}
\renewcommand\headrulewidth{1pt}
\fancyhead[L]{TP 4}
\fancyhead[R]{Saux Paulhenry}
\title{Compte rendu de TP 4}
\date{}
\author{}

\geometry{a4paper}
% ----------- Création des commandes de couleur ----------
\makeatletter
\newcommand{\coloruline}[2]{%
        \UL@protected\def\temp@uline{\leavevmode \bgroup
            \markoverwith{\textcolor{#1}{\rule[-0.5ex]{2pt}{0.4pt}}}%
        \ULon}%
        \temp@uline{#2}%
    }

    \newcommand{\coloruuline}[2]{%
        \UL@protected\def\temp@uuline{\leavevmode \bgroup
            \UL@setULdepth
            \ifx\UL@on\UL@onin \advance\ULdepth2.8\p@\fi
            \markoverwith{\textcolor{#1}{\lower\ULdepth\hbox
                {\kern-.03em\vbox{\hrule width.2em\kern1\p@\hrule}\kern-.03em}}}%
        \ULon}%
        \temp@uuline{#2}%
    }

\makeatother

\renewcommand{\thesection}{\arabic{section}}
\renewcommand{\thesubsection}{\arabic{subsection}}
\renewcommand{\thesubsubsection}{\arabic{subsubsection}}
\begin{document}
% ---------------- Modification des titres de niveau 1,2 et 3 --------------------
\titleformat
{\section} % command
%[display] % shape
{\Large} % format
{\coloruuline{red}{\thesection. }} % label
{0ex} % sep
{\coloruuline{red}} % before-code
[] % after-code

\titleformat
{\paragraph} % command
%[display] % shape
{\Large} % format
{\coloruline{NavyBlue}{\theparagraph. }} % label
{0ex} % sep
{\coloruline{NavyBlue}} % before-code
[] % after-code

\titleformat
{\subsection} % command
%[display] % shape
{\Large} % format
{\coloruline{red}{\thesection.\thesubsection. }} % label
{0ex} % sep
{\coloruline{red}} % before-code
[] % after-code

\titleformat
{\subsubsection} % command
%[display] % shape
{\large} % format
{\coloruline{ForestGreen}{\thesection.\thesubsection.\thesubsubsection. }} % label
{0ex} % sep
{\coloruline{ForestGreen}} % before-code
[] % after-code

\titleformat
{\chapter} % command
%[display] % shape
{\large} % format
{\coloruline{NavyBlue}{\thechapter \:- }} % label
{0ex} % sep
{\coloruline{NavyBlue}} % before-code
[] % after-code


% ---------------- Début du corps du document ------------------------------------

\section{Introduction}
\subsection{Objectifs}
L'objectif est de mesurer le rapport \(e/m_e\) de deux façons différentes afin de déterminer quel protocole est le meilleur. Il s'agira dans les deux cas d'utiliser le mouvement d'électrons dans un champ magnétique uniforme.
\subsection{Sommaire}
Ce TP se décompose en 2 parties :
\begin{enumerate}
 \item Mesure en maintenant le rayon fixe (Protocole 1)
 \item Mesure en maintenant la tension fixe  (Protocole 2)
\end{enumerate}
\section{Protocole}
\subsection{Matériel}
\begin{itemize}
\item Tube à faisceau électronique filiforme
\item 2 Bobines en configuration Helmholtz
\item Dispositif de Mesure
\item 2 Multimètres
\item Générateurs respectant les conditions :
\begin{itemize}
 \item 6.3V, 1A
 \item 0-500V, 50 mA
 \item 0-50V, 10mA
 \item 0-3A
\end{itemize}
\item Cables
\end{itemize}
\subsection{Description du montage expérimental}
Nous branchons les éléments du circuit électrique comme indiqué sur la figure ci-dessous. Le voltmètre mesure la tension entre l'anode et la cathode du Tube à faisceau électronique filiforme et l'ampèremètre l'intensité fournie aux bobines.
\begin{figure}
 \begin{center}
\includegraphics[scale=1.2]{schema}
\end{center}
\caption{Schema du montage expérimental}
\end{figure}
\subsection{Mesures à effectuer}
Nous procédons selon 2 protocoles.
\subsubsection{Protocole 1 : Rayon constant}
Nous maintenons le rayon constant durant tout le protocole. On choisi un rayon le plus grand possible, tout en gardant de bonnes conditions de lisibilité. Nous avons choisi de prendre un rayon de 11.6 cm.
\begin{enumerate}
 \item Placer la tension du générateur à 130V environ
 \item Placer le premier marqueur de mesure à l'origine du faisceau
 \item Réglier l'intensité des bobines pour observer un diamètre de 11.6cm environ
 \item Relever la paire Tension/Intensité
 \item Augmenter la tension du générateur jusqu'à 300V par pas de 20V et refaire les étapes 2 et 3
\end{enumerate}
\subsubsection{Protocole 2 : Tension constante}
\begin{enumerate}
 \item Placer la tension du générateur à 200V environ
  \item Placer le premier marqueur de mesure à l'origine du faisceau
 \item Réglier l'intensité des bobines à 1A
 \item Relever le diamètre du cercle de la trajectoire adoptée par les électrons
 \item Pour des valeurs de I de 1 à 2.1A par pas de 0.1, refaire l'étape précédente.q
\end{enumerate}
\section{Mesures}
\subsection{Données}
\begin{table}[H]
\begin{center}
\begin{tabular}{ll}
U (V) & I (A) \\
$1,31 \times 10^2$ & $5,83 \times 10^{-1}$ \\
$1,51 \times 10^2$ & $8,07 \times 10^{-1}$ \\
$1,70 \times 10^2$ & $8,68 \times 10^{-1}$ \\
$1,91 \times 10^2$ & $9,80 \times 10^{-1}$ \\
$2,09 \times 10^2$ & $1,06 \times 10^{0}$ \\
$2,31 \times 10^2$ & $1,07 \times 10^{0}$ \\
$2,49 \times 10^2$ & $1,17 \times 10^{0}$ \\
$2,70 \times 10^2$ & $1,23 \times 10^{0}$ \\
$2,89 \times 10^2$ & $1,32 \times 10^{0}$ \\
$3,10 \times 10^2$ & $1,41 \times 10^{0}$
\end{tabular}
\end{center}
\caption{Intensité et Tensions mesurés dans le cadre du protocole 2}
\end{table}

\begin{table}[H]
 \begin{center}
\begin{tabular}{ll}
I (A) & D (m)\\
$1,00 \times 10^0$ & $1,22 \times 10^{-1}$\\
$1,09 \times 10^0$ & $1,13 \times 10^{-1}$\\
$1,21 \times 10^0$ & $1,05 \times 10^{-1}$\\
$1,31 \times 10^0$ & $9,90 \times 10^{-2}$\\
$1,40 \times 10^0$ & $8,90 \times 10^{-2}$\\
$1,51 \times 10^0$ & $8,20 \times 10^{-2}$\\
$1,61 \times 10^0$ & $7,70 \times 10^{-2}$\\
$1,70 \times 10^0$ & $7,30 \times 10^{-2}$\\
$1,81 \times 10^0$ & $6,90 \times 10^{-2}$\\
$1,90 \times 10^0$ & $6,40 \times 10^{-2}$\\
$2,01 \times 10^0$ & $6,10 \times 10^{-2}$\\
$2,11 \times 10^0$ & $5,80 \times 10^{-2}$
\end{tabular}
\end{center}
\caption{Intensité et diamètres mesurés dans le cadre du protocole 2}
\end{table}


\subsection{Incertitudes}
\subsubsection{Protocole 1}
Les incertitudes de la tension et de l'intensité sont liées uniquement à la précision des instruments de mesure. Ces dernières sont négligeables devant l'incertitude systématique liée au rayon de la trajectoire des électrons, qui est systématique et estimée à 7mm. Cette valeur correspond à l'épaisseur du faisceau. En effet, ce dernier n'est pas d'une bonne qualité et est moins concentré. Cela empeche de déterminer avec précision la zone où la trajectoire a lieu (on a un gradient au lieu d'une coupure nette).
\subsubsection{Protocole 2}
Les incertitudes de la tension et de l'intensité sont liées uniquement à la précision des instruments de mesure. Ces dernières sont négligeables devant l'incertitude liée à la masure du diamètre. En effet, cette dernière  est réalisée par le mesure entre 2 repères amovibles, l'un à l'origine du faiseacu et l'autre diamétralement opposé. Il y a dès lors une incertitude systématique liée à la position du premier marqueur, estimée à \(\pm 1\)mm et celle liée à la position du second marqueur. Cette dernière doit prendre en compte les différences de parallaxe lors des obervsations et la largeur du faisceau, qui peut empecher de connaitre avec précision la trajectoire. Dans les conditions de ce protocole, le fasceau produit était de meilleure qualité que lors du protocole 1, avec une épaisseur très faible, d'environ \(2\)mm. On estimera les incertitudes liées à la parralaxe à \(2\)mm. Enfin, comme les marqueurs ne bougent pas de façon perpendiculaire aux axes (la pertie haute peut avoir un décalage de 1 mm avec la partie du bas), on ajoute une incertitude de 2mm au résultat, bien qu'une attention particulière ai étée portée à ce problème dans la manipulation. L'incertitude liée au diamètre est donc de \(2+2+1=5mm\).
\section{Graphiques}
\begin{figure}[H]
\begin{center}
  \includegraphics[scale=0.5]{1}
\end{center}
\caption{Évolution du champ magnétique (qui dépend de I) en fonction de U}
\end{figure}
\begin{figure}[H]
\begin{center}
  \includegraphics[scale=0.5]{2}
\end{center}
\caption{Évolution du rayon au carré en fonction de \(\frac{1}{B^2}\)}
\end{figure}
\section{Exploitation des résultats}
On se servira dans les 2 modélisations des relations \[B^2 = U\cdot (\frac{2m_e}{e r^2})\] et \[B^2 = \frac{64}{125}(\frac{\mu_0NI}{R})^2 \] avec \(I\) l'itensité mésurée pour atteindre le rayon fixe, N le nombre de spires des bobines, R le rayon des bobines et \(\mu_0\) une constante.
\subsection{Protocole 1}
On peut déduire la valeur du rapport en traçant \(B^2\) en fonction de \(U\) et en modélisant par une fonction linéaire de pente A : \[B^2 = f(U) = A\cdot U\]

On peut alors déduire la valeur recherchée et on aura :

\[u(e/m_e) = \frac{e}{m_e}[4(\frac{u(r)}{r})^2+(\frac{u(A)}{A})^2]^{1/2}\]

On obtient \(\frac{e}{m_e} = (177,7\pm8,5)\cdot 10^9S.I\). L'encadrement inclue bien la valeur théorique.

L'erreur relative par rapport à la valeur théorique ($1,75882 \cdot 10^{11} SI$) est de 1,03\%.
\subsection{Protocole 2}
On déduit des m\^emes relations que précédement que pour obtenir la valeur recherchée, on peut tracer \(r^2\) en fonction de \(B^{-2}\) et modéliser par une fonction linéaire dont la pente C sera \(2U\frac{m_e}{e}\).

En effet, on peut écrire \(r^2=2UB^{-2}A^{-1}\)

$B^{-2}$ est calculé avec \[B^{-2} = \big(\frac{64}{125}(\frac{\mu_0NI}{R})^2\big)^{-1}\]

On trace \[r^2 = f(B^{-2}) = C\cdot B^{-2}\]

On peut alors déduire le rapport de cette pente et on aura :
\[u(e/m_e) = \frac{e}{m_e}[(\frac{u(U)}{U})^2+(\frac{u(C)}{C})^2]^{1/2}\]

On obtient de cette façon que \((\frac{e}{m_e} = 173,4 \pm 2,1)\cdot 10^9\) SI.

L'encadrement inclue bien la valeur théorique. De plus, l'erreur relative par rapport à la valeur théorique est de 1,41\%.
\section{Conclusion}
On remarque que, bien que le protocole 1 introduise moins d'incertitude sur les mesures, en raison de la mauvaise qualité du faisceau utilisé lors de cette partie, c'est le protocole 2 qui donne un résultat avec l'incertitude la plus faible. Ainsi, la façon de mesurer le diamètre dans le premier protocole ne parvient pas à compenser la qualité du faisceau.

On note cependant que bien que le Protocole 1 ait une incertitude plus grande que le Protocole 2, il a donné la valeur qui est finalement plus proche de la valeur théorique (erreur relative de 1,03\% contre 1,41\% pour le Protocole 2). Dans les 2 cas, les valeurs obtenues sont très satisfaisantes avec une erreur relative de l'ordre de \(1\%\). L'utilisation d'un faisceau électronique filiforme est donc une méthode jugée satisfaisante pour déterminer la valeur du rapport \(\frac{e}{e_m}\).
\end{document}
